% GENERATED FILE. Edit canonical lessons, then run npm run generate:books.
% canonical-source-hash: fnv1a64:bec19c3235bb6a45
% canonical-lessons: LA-C299-narratio, LA-C299-commemoro, LA-C299-memoro, LA-C299-refero, LA-C299-inquam

\chapter{Nārrātiō}
\label{ch:la-narratio}
\hlchaptermodality{299}

\begin{chapteropening}
\textbf{By the end of this chapter:} \emph{I can talk about telling a story: a narrative, to recall, to relate, to report and to quote words.}

Complete the last lesson of chapter 299: I can talk about telling a story: a narrative, to recall, to relate, to report and to quote words.
\end{chapteropening}

\section[nārrātiō, nārrātiōnis]{nārrātiō, nārrātiōnis — a narrative, a telling}

\label{lesson:LA-C299-narratio}



\begin{quote}
Before the new one: say the Latin for to come along, then the Latin for next.
\end{quote}

\subsection*{You'll want to know: nārrātiō, nārrātiōnis}
\textbf{nārrātiō, nārrātiōnis} — \textquotedblleft{}a narrative, a telling\textquotedblright{}.

\begin{quote}
Nārrātiō senis longa sed iūcunda erat.
\end{quote}

\begin{quote}
\emph{The old man's story was long but pleasant.}
\end{quote}

Feminine, 3rd declension, from \emph{nārrō} (taught). In a Roman speech, the \emph{nārrātiō} set out the facts. English \emph{narration}.

\begin{cousinweb}
\textbf{nārrātiō} is built on \emph{nārrō}, \textquotedblleft{}to tell, to relate\textquotedblright{} (EtymDB), which Lewis \& Short link with \emph{gnārus}, \textquotedblleft{}knowing\textquotedblright{}: making known. English \emph{narration} comes from it.
\end{cousinweb}

\begin{grammarlens}[title={What you've built}]
One more word for telling, explaining and describing at B1.
\end{grammarlens}

\subsection*{Your turn}
\begin{itemize}
\raggedright
  \item \emph{Say it:} \emph{nārrātiō, nārrātiōnis}
  \item \emph{Say it:} \emph{nārrātiō, nārrātiōnis}, once more
  \item \emph{Say it:} a sentence of your own with \emph{nārrātiō, nārrātiōnis}
  \item \emph{Recall:} say the Latin for to come along, then the Latin for next, then say \emph{nārrātiō, nārrātiōnis} again
\end{itemize}

\begin{tcolorbox}[breakable,colback=teal!4,colframe=teal!35!black,title={Before you move on}]
What does nārrātiō, nārrātiōnis mean? (\textquotedblleft{}A narrative, a telling\textquotedblright{}.) Say it once more.
\end{tcolorbox}
\section[commemorō, commemorāre]{commemorō, commemorāre — to recall, to mention}

\label{lesson:LA-C299-commemoro}



\begin{quote}
Before the new one: say the Latin for next, then the Latin for a narrative.
\end{quote}

\subsection*{You'll want to know: commemorō, commemorāre}
\textbf{commemorō, commemorāre} — \textquotedblleft{}to recall, to mention\textquotedblright{}.

\begin{quote}
Senex victōriās veterēs commemorābat.
\end{quote}

\begin{quote}
\emph{The old man was recalling old victories.}
\end{quote}

\emph{commemorō, commemorāre, commemorāvī, commemorātum}: to bring something fully to mind, and so to mention it. English \emph{commemorate}.

\begin{cousinweb}
\textbf{commemorō} is built on \emph{memorō}, \textquotedblleft{}to bring to mind, to tell of\textquotedblright{}, with \emph{con-} (EtymDB); \emph{memorō} comes from \emph{memor}, \textquotedblleft{}mindful\textquotedblright{}. English \emph{commemorate} comes from it.
\end{cousinweb}

\begin{grammarlens}[title={What you've built}]
One more word for telling, explaining and describing at B1.
\end{grammarlens}

\subsection*{Your turn}
\begin{itemize}
\raggedright
  \item \emph{Say it:} \emph{commemorō, commemorāre}
  \item \emph{Say it:} \emph{commemorō, commemorāre}, once more
  \item \emph{Say it:} a sentence of your own with \emph{commemorō, commemorāre}
  \item \emph{Recall:} say the Latin for next, then the Latin for a narrative, then say \emph{commemorō, commemorāre} again
\end{itemize}

\begin{tcolorbox}[breakable,colback=teal!4,colframe=teal!35!black,title={Before you move on}]
What does commemorō, commemorāre mean? (\textquotedblleft{}To recall, to mention\textquotedblright{}.) Say it once more.
\end{tcolorbox}
\section[memorō, memorāre]{memorō, memorāre — to tell of, to relate}

\label{lesson:LA-C299-memoro}



\begin{quote}
Before the new one: say the Latin for a narrative, then the Latin for to recall.
\end{quote}

\subsection*{You'll want to know: memorō, memorāre}
\textbf{memorō, memorāre} — \textquotedblleft{}to tell of, to relate\textquotedblright{}.

\begin{quote}
Poētae multa dē bellīs memorant.
\end{quote}

\begin{quote}
\emph{Poets tell many things about wars.}
\end{quote}

\emph{memorō, memorāre, memorāvī, memorātum}. Poetic and historical in tone; Vergil opens his poem by asking the Muse to tell him the causes.

\begin{cousinweb}
\textbf{memorō} is built on \emph{memor}, \textquotedblleft{}mindful\textquotedblright{} (EtymDB), which is inherited from a Proto-Indo-European root meaning \textquotedblleft{}to remember\textquotedblright{}. English \emph{memory} comes from \emph{memoria}, of the same family.
\end{cousinweb}

\begin{grammarlens}[title={What you've built}]
One more word for telling, explaining and describing at B1.
\end{grammarlens}

\subsection*{Your turn}
\begin{itemize}
\raggedright
  \item \emph{Say it:} \emph{memorō, memorāre}
  \item \emph{Say it:} \emph{memorō, memorāre}, once more
  \item \emph{Say it:} a sentence of your own with \emph{memorō, memorāre}
  \item \emph{Recall:} say the Latin for a narrative, then the Latin for to recall, then say \emph{memorō, memorāre} again
\end{itemize}

\begin{tcolorbox}[breakable,colback=teal!4,colframe=teal!35!black,title={Before you move on}]
What does memorō, memorāre mean? (\textquotedblleft{}To tell of, to relate\textquotedblright{}.) Say it once more.
\end{tcolorbox}
\section[referō, referre]{referō, referre — to report; to carry back}

\label{lesson:LA-C299-refero}



\begin{quote}
Before the new one: say the Latin for to recall, then the Latin for to tell of.
\end{quote}

\subsection*{You'll want to know: referō, referre}
\textbf{referō, referre} — \textquotedblleft{}to report; to carry back\textquotedblright{}.

\begin{quote}
Nūntius omnia quae vīderat ducī rettulit.
\end{quote}

\begin{quote}
\emph{The messenger reported to the leader all that he had seen.}
\end{quote}

\emph{referō, referre, rettulī, relātum}, irregular like \emph{ferō}. \emph{Ad senātum referre}: to put a matter before the senate.

\begin{cousinweb}
\textbf{referō} is built on \emph{ferō}, \textquotedblleft{}to carry\textquotedblright{}, with \emph{re-}, \textquotedblleft{}back\textquotedblright{} (EtymDB): to carry back, and so to bring back word. English \emph{refer} comes from it, and \emph{relate} from its participle.
\end{cousinweb}

\begin{grammarlens}[title={What you've built}]
One more word for telling, explaining and describing at B1.
\end{grammarlens}

\subsection*{Your turn}
\begin{itemize}
\raggedright
  \item \emph{Say it:} \emph{referō, referre}
  \item \emph{Say it:} \emph{referō, referre}, once more
  \item \emph{Say it:} a sentence of your own with \emph{referō, referre}
  \item \emph{Recall:} say the Latin for to recall, then the Latin for to tell of, then say \emph{referō, referre} again
\end{itemize}

\begin{tcolorbox}[breakable,colback=teal!4,colframe=teal!35!black,title={Before you move on}]
What does referō, referre mean? (\textquotedblleft{}To report; to carry back\textquotedblright{}.) Say it once more.
\end{tcolorbox}
\section[inquam]{inquam — I say (quoting words)}

\label{lesson:LA-C299-inquam}



\begin{quote}
Before the new one: say the Latin for to tell of, then the Latin for to report.
\end{quote}

\subsection*{You'll want to know: inquam}
\textbf{inquam} — \textquotedblleft{}I say (quoting words)\textquotedblright{}.

\begin{quote}
'Venī mēcum,' inquit pater, 'et omnia vidēbis.'
\end{quote}

\begin{quote}
\emph{'Come with me,' said father, 'and you will see everything.'}
\end{quote}

Defective: mostly \emph{inquam, inquis, inquit}, set inside the words quoted. \emph{Inquit} means 'he says' or 'he said'.

\begin{cousinweb}
\textbf{inquam} is an inherited verb that EtymDB traces to a Proto-Indo-European form \emph{*ensk\textsuperscript{w}e-}, \textquotedblleft{}to say\textquotedblright{}. Only a few forms are in use (Lewis \& Short).
\end{cousinweb}

\begin{grammarlens}[title={What you've built}]
One more word for telling, explaining and describing at B1.
\end{grammarlens}

\subsection*{Your turn}
\begin{itemize}
\raggedright
  \item \emph{Say it:} \emph{inquam}
  \item \emph{Say it:} \emph{inquam}, once more
  \item \emph{Say it:} a sentence of your own with \emph{inquam}
  \item \emph{Recall:} say the Latin for to tell of, then the Latin for to report, then say \emph{inquam} again
\end{itemize}

\begin{tcolorbox}[breakable,colback=teal!4,colframe=teal!35!black,title={Before you move on}]
What does inquam mean? (\textquotedblleft{}I say (quoting words)\textquotedblright{}.) Say it once more.
\end{tcolorbox}
